% Session 3 -- two ways to gap the honeycomb Dirac cones, and why only one of
% them gives a Chern insulator.
% Placement: sec. "Gapping the Dirac cone", after equation (3.8).
%
% LAYOUT NOTE: pure tikzpicture, bare canvas coordinates throughout. The only
% local style is a tiny-label style, declared inside this picture.
\begin{aqmfigure}{fig:s03:masses}{The full content of Section 3.10 as one counting of Chern numbers.
Equation (3.8) gives each gapped Dirac point a Chern number $\tfrac12\,\mathrm{sgn}(m)\,
\mathrm{sgn}(\lambda)$, and the two Dirac points of the honeycomb lattice are time-reversal partners
with opposite chirality, $\lambda=+1$ at $K$ and $\lambda=-1$ at $K'$. So the result depends only on
how the mass term itself behaves under time reversal. (a) A staggered sublattice potential is an
ordinary scalar potential, which is even under time reversal, so $m$ has the same sign at both
valleys. The two half-integers are then opposite and they cancel, which gives an ordinary insulator
with the same kind of gap as boron nitride. (b) The Haldane coupling is a pattern of circulating
orbital currents, which is odd under time reversal, so $m$ changes sign between the valleys. The two
half-integers are then equal and they add, which gives $C=\pm1$. This is a real Chern insulator with
quantized $\sigma_{xy}$ and chiral edge states, built only from lattice hoppings, with no external
magnetic field and zero net flux. The two gaps look the same in the bulk spectrum, and the difference
appears only in the wavefunctions.}

\begin{tikzpicture}[x=1cm, y=1cm]
\tikzset{aqmlabtiny/.style={font=\tiny, inner sep=0.8pt}}

\def\hexR{0.66}

% =====================================================================
% header: the Brillouin zone and the rule
% =====================================================================
\begin{scope}[shift={(1.00,0)}]
  \draw[aqm, draw=black] (0:0.86) \foreach \a in {60,120,...,360} {-- (\a:0.86)} -- cycle;
  \foreach \a in {0,120,240} {
    \node[circle, fill=aqmaccent, draw=aqmaccent, minimum size=4.5pt, inner sep=0pt]
      at (\a:0.86) {};
  }
  \foreach \a in {60,180,300} {
    \node[circle, fill=white, draw=black, minimum size=4.5pt, inner sep=0pt]
      at (\a:0.86) {};
  }
  \node[aqmlabsmall, anchor=north west] at (0:0.94)   {$K$};
  \node[aqmlabsmall, anchor=south west] at (60:0.94)  {$K'$};
  \node[aqmlabsmall, anchor=north] at (0,-1.06) {Brillouin zone};
\end{scope}

\node[aqmlab, anchor=west, align=left, text width=9.9cm] at (2.75,0)
  {\textbf{The rule (3.8):}\quad
   $C_{\text{Dirac}} = \tfrac12\,\mathrm{sgn}(m)\,\mathrm{sgn}(\lambda)$.\\[3pt]
   The lattice fixes the chirality: $\lambda = +1$ at $K$ and $\lambda = -1$ at $K'$,\\[3pt]
   because $K$ and $K'$ are time-reversal partners; only the \emph{sum}\\[3pt]
   over the two valleys has to be an integer.};

\draw[aqm, draw=aqmmuted, line width=\aqmthinwidth] (-0.35,-1.72) -- (12.7,-1.72);

% =====================================================================
% (a) staggered sublattice potential -> C = 0
% =====================================================================
\begin{scope}[shift={(0,-3.55)}]

\node[aqmpanel, anchor=west] at (-0.35,1.45) {(a) staggered sublattice potential};

% -- lattice icon --
\begin{scope}[shift={(1.05,-0.12)}]
  \draw[aqmbond] (0:\hexR) \foreach \a in {60,120,...,360} {-- (\a:\hexR)} -- cycle;
  \foreach \a in {0,120,240} {
    \draw[aqmbond] (\a:\hexR) -- (\a:{\hexR+0.30});
    \node[aqmsite] at (\a:\hexR) {};
    \node[aqmlabtiny] at (\a:{\hexR+0.46}) {$+\Delta$};
  }
  \foreach \a in {60,180,300} {
    \draw[aqmbond] (\a:\hexR) -- (\a:{\hexR+0.30});
    \node[aqmsiteopen] at (\a:\hexR) {};
    \node[aqmlabtiny] at (\a:{\hexR+0.46}) {$-\Delta$};
  }
\end{scope}

\node[aqmlab, anchor=west, align=left, text width=4.55cm] at (2.62,-0.12)
  {a scalar on-site potential:\\[1pt]
   \textbf{even} under time reversal,\\[1pt]
   so $m$ has the \textbf{same} sign\\[1pt]
   at $K$ and $K'$.};

\node[aqmboxaccent, anchor=west, align=left, text width=4.7cm] at (7.45,-0.12)
  {$C_K \;=\; \tfrac12\,(+)(+) \;=\; +\tfrac12$\\[2pt]
   $C_{K'} \;=\; \tfrac12\,(+)(-) \;=\; -\tfrac12$\\[3pt]
   \hrule height 0.4pt \vspace{3pt}
   $C = 0$: the two cancel.\\[1pt]
   An \textbf{ordinary insulator}.};

\end{scope}

% =====================================================================
% (b) Haldane coupling -> C = +-1
% =====================================================================
\begin{scope}[shift={(0,-7.15)}]

\node[aqmpanel, anchor=west] at (-0.35,1.45) {(b) Haldane coupling};

% -- lattice icon --
\begin{scope}[shift={(1.05,-0.12)}]
  \draw[aqmbond] (0:\hexR) \foreach \a in {60,120,...,360} {-- (\a:\hexR)} -- cycle;
  \foreach \a in {0,120,240} {
    \draw[aqmbond] (\a:\hexR) -- (\a:{\hexR+0.30});
    \node[aqmsite] at (\a:\hexR) {};
  }
  \foreach \a in {60,180,300} {
    \draw[aqmbond] (\a:\hexR) -- (\a:{\hexR+0.30});
    \node[aqmsiteopen] at (\a:\hexR) {};
  }
  % circulating orbital current inside the hexagon
  \draw[aqm, draw=aqmbroken, line width=1.1pt, -{Stealth[length=5pt,width=4pt]}]
    (-95:0.34) arc (-95:215:0.34);
  \node[aqmlabtiny, text=aqmbroken, anchor=north] at (0,-0.86) {loop current};
\end{scope}

\node[aqmlab, anchor=west, align=left, text width=4.55cm] at (2.62,-0.12)
  {circulating orbital currents:\\[1pt]
   \textbf{odd} under time reversal,\\[1pt]
   so $m$ \textbf{flips} sign between\\[1pt]
   $K$ and $K'$.};

\node[aqmboxbroken, anchor=west, align=left, text width=4.7cm] at (7.45,-0.12)
  {$C_K \;=\; \tfrac12\,(+)(+) \;=\; +\tfrac12$\\[2pt]
   $C_{K'} \;=\; \tfrac12\,(-)(-) \;=\; +\tfrac12$\\[3pt]
   \hrule height 0.4pt \vspace{3pt}
   $C = +1$: the two add.\\[1pt]
   A \textbf{Chern insulator}, no field.\\[1pt]
   (Reversing the coupling gives $-1$.)};

\end{scope}

\end{tikzpicture}

\end{aqmfigure}
